\documentclass[a4paper,11pt,fleqn]{article}%fleqn pour aligner les equations à gauche
\usepackage{../../Styles/Cours}

\newcommand{\chnum}{Chapitre 11}
\newcommand{\chtitle}{Mouvement circulaire uniforme}
\newcommand{\dtype}{TP}


\begin{document}
\maketitle

\section{Le MCU}
Le mouvement d'un objet est analysé deux fois, chacune avec un intervalle de temps différent.
\begin{table}[h!]
	\begin{tabular}{|>{\centering}m{9cm} | >{\centering \arraybackslash}m{9cm}|}
	\hline
	Intervalle de temps $\Delta t$ & Intervalle de temps $\dfrac{\Delta t}{3}$\\
	\hline
	\begin{tikzpicture}
		\draw[blue] (0,0) circle (4);
		\foreach \k in {0,45,90,135,180,315}{
			\pgfmathsetmacro{\x}{4*cos(\k)}
			\pgfmathsetmacro{\y}{4*sin(\k)}
			%\node at (\x,\y){\Huge +};
			\draw (\x,\y) node{\Huge + };
		}
		\node[above] (Mi) at (0,4){\Huge $M_i$};
		\draw[->,thick,>=latex, out=90, in =220] (-4.5,0) to (-3.19,3.19);
	\end{tikzpicture}&
	\begin{tikzpicture}
		\draw[blue] (0,0) circle (4);
		\foreach \k in {0,15,30,45,60,75,90,105,120,135,150,165,180,315,330,345}{
			\pgfmathsetmacro{\x}{4*cos(\k)}
			\pgfmathsetmacro{\y}{4*sin(\k)}
			%\node at (\x,\y){\Huge +};
			\draw (\x,\y) node{\Huge + };
		}
		\node[above] (Mi) at (0,4){\Huge $M_i$};
		\draw[->,thick,>=latex, out=90, in =220] (-4.5,0) to (-3.19,3.19);
	\end{tikzpicture}\\
	\hline
	\end{tabular}
\end{table}
\begin{enumerate}
	\item Décrire le mouvement de l'objet.\vspace{2cm}
	\item Représenter les vecteurs vitesse $\vec{v_i}$ et $\vv{v_{i+1}}$ dans les deux situations, sans considération d'échelle, en bleu.\vspace{2cm}
	\item Construire le vecteur variation de vitesse $\Delta \vec{v_i}$, en rouge.\vspace{2cm}
	\item En déduire la direction et le sens du vecteur accélération $\vec{a_i}$ lors d'un mouvement circulaire uniforme pour un intervalle de temps infiniment court.\vspace{2cm}
\end{enumerate}
\newpage
\section{Accélération dans le repère de Frenet}
\begin{enumerate}
	\item Construire, sur le document fourni, les vecteurs accélération en $M_2$ et $M_9$, en rouge.\vspace{2cm}
	\item En déduire, par lecture graphique, les coordonnées $a_T$ et $a_N$ du vecteur accélération dans le repère de Frenet en $M_2$ et $M_9$. \vspace{2cm}
	\item Montrer que ces coordonnées sont en accord avec l'expression suivante :
	$$	\vec{a}=\dfrac{dv}{dt} \cdot \vec{T} + \dfrac{v^2}{R} \cdot \vec{N} $$
\end{enumerate}
\begin{figure}[b!]
	\centering
\begin{tikzpicture}
	\draw[blue] (0,0) circle (4);
	\foreach \k in {1,...,26}{
		\pgfmathsetmacro{\i}{\k*360/26}
		\pgfmathsetmacro{\x}{4*cos(\i)}
		\pgfmathsetmacro{\y}{4*sin(\i)}
		%\node at (\x,\y){\Huge +};
		\draw (\x,\y) node{\Huge +};
	}
	\foreach \l in {-4,...,4}{
		\draw[line width=0.05cm](\l,-4.7)--(\l,-4.3);
		\draw[line width=0.05cm](-4.7,\l)--(-4.3,\l);
		\node[below] at (\l,-4.7){\l};
		\node[left] at (-4.7,\l){\l};
	}
	\node[above] (M2) at (0.5,4.2){\huge $M_2$};
	\node[right] (M9) at (-4,0){\huge $M_9$};
	%\draw[->,thick,>=latex, out=90, in =220] (-4.5,0) to (-3.19,3.19);
	\draw[-Latex, line width=.05cm] (-4.5,-4.5) to (4.5,-4.5);
	\draw[-Latex, line width=.05cm] (-4.5,-4.5) to (-4.5,4.5);
	\draw[line width=.05cm] (-4.5,0)--(4.5,0);
	\draw[line width=.05cm] (0,-4.5)--(0,4.5);
	\draw[color=gray] (-4.5,-4.5) grid(4.5,4.5);

	\node[right] at(4.5,-4.5){\huge x(m)};
	\node[above] at(-4.5,4.5){\huge y(m)};
\end{tikzpicture}\\
Durée entre deux positions consécutives $\Delta t$ = \SI{25}{ms}
\end{figure}
\end{document}
